%% slides/beamer-rpd.tex ------------------------------------------------------------
%% Flat beamer theme for the RPD conference talk.
%%
%% Loaded through `include-in-header:` in rpd-talk.qmd. Pandoc appends this file to
%% the `header-includes` variable, which Quarto's beamer template expands at
%% template.tex:90 -- i.e. AFTER \usetheme (l.72) and AFTER the package named by
%% `fontfamily:` (l.89), and BEFORE \begin{document} (l.99). That ordering is the
%% whole reason the plain \setbeamertemplate overrides below win without needing
%% \AtBeginDocument.
%%
%% Dependencies: beamer, xcolor, graphicx -- all in collection-latexrecommended,
%% already installed by .devcontainer/Dockerfile. Nothing else is used, because the
%% deck must render under `docker run --network none`. In particular metropolis,
%% pgfopts, fira, adjustbox, tcolorbox and appendixnumberbeamer are NOT in the image
%% and are deliberately worked around rather than required.
\makeatletter

%% ---- palette -----------------------------------------------------------------------
%% Near-black ink, one cool accent for structure, one warm accent for emphasis. Both
%% accents clear 4.5:1 against white and stay distinguishable under deuteranopia.
%% R/slides_en.R defines the same hex values, so figures and chrome match exactly.
\definecolor{rpdink}{HTML}{1A2028}
\definecolor{rpdmute}{HTML}{6E7884}
\definecolor{rpdrule}{HTML}{DADEE2}
\definecolor{rpdaccent}{HTML}{00506B}
\definecolor{rpdalert}{HTML}{B4532A}

\setbeamercolor{normal text}{fg=rpdink, bg=white}
\setbeamercolor{structure}{fg=rpdaccent}
\setbeamercolor{alerted text}{fg=rpdalert}
\setbeamercolor{frametitle}{fg=rpdink, bg=}      % empty bg => no title band
\setbeamercolor{framesubtitle}{fg=rpdmute, bg=}
\setbeamercolor{title}{fg=rpdink}
\setbeamercolor{subtitle}{fg=rpdaccent}
\setbeamercolor{author}{fg=rpdink}
\setbeamercolor{institute}{fg=rpdmute}
\setbeamercolor{date}{fg=rpdmute}
\setbeamercolor{section page}{fg=rpdink}
\setbeamercolor{itemize item}{fg=rpdaccent}
\setbeamercolor{itemize subitem}{fg=rpdmute}
\setbeamercolor{enumerate item}{fg=rpdaccent}
\setbeamercolor{block title}{fg=white, bg=rpdaccent}
\setbeamercolor{block body}{fg=rpdink, bg=rpdaccent!7}
\setbeamercolor{caption name}{fg=rpdmute}
\setbeamercolor{footline}{fg=rpdmute}
\setbeamercolor{footnote}{fg=rpdmute}
\setbeamercolor{bibliography entry author}{fg=rpdink}
\setbeamercolor{bibliography entry title}{fg=rpdink}
\setbeamercolor{bibliography entry location}{fg=rpdmute}
\setbeamercolor{bibliography entry note}{fg=rpdmute}

%% ---- fonts -------------------------------------------------------------------------
%% `fontfamily: libertinus` has already loaded libertinus-otf (libertinus.sty branches
%% on \iftutex), so \sffamily is Libertinus Sans, \rmfamily is Libertinus Serif, and
%% unicode-math points at LibertinusMath-Regular.otf. Do NOT use `mainfont:` in the
%% YAML instead: Libertinus is invisible to fontconfig -- it lives only in the texmf
%% tree -- so \setmainfont would fail to find it, and setting it also triggers
%% \usefonttheme{serif} at template.tex:81, turning the whole deck serif.
%% `professionalfonts` stops beamer from overriding the math setup unicode-math just made.
\usefonttheme{professionalfonts}
\setbeamerfont{title}{family=\sffamily, size=\LARGE, series=\bfseries}
\setbeamerfont{subtitle}{family=\sffamily, size=\large, series=\mdseries}
\setbeamerfont{author}{family=\sffamily, size=\normalsize}
\setbeamerfont{institute}{family=\sffamily, size=\small}
\setbeamerfont{date}{family=\sffamily, size=\small}
\setbeamerfont{frametitle}{family=\sffamily, size=\large, series=\bfseries}
\setbeamerfont{framesubtitle}{family=\sffamily, size=\small, series=\mdseries}
\setbeamerfont{section page}{family=\sffamily, size=\LARGE, series=\bfseries}
\setbeamerfont{block title}{family=\sffamily, size=\normalsize, series=\bfseries}
\setbeamerfont{footline}{family=\sffamily, size=\scriptsize}
\setbeamerfont{caption}{size=\scriptsize}
\setbeamerfont{caption name}{size=\scriptsize, series=\mdseries}
\setbeamerfont{footnote}{size=\scriptsize}

%% ---- rhythm ------------------------------------------------------------------------
\setbeamersize{text margin left=0.9cm, text margin right=0.9cm}
\setlength{\parskip}{0.55em}          % parskip.sty is already loaded by common.latex
\renewcommand{\arraystretch}{1.08}    % breathing room for dense kableExtra tables

%% ---- chrome ------------------------------------------------------------------------
\setbeamertemplate{navigation symbols}{}
\setbeamertemplate{sidebar right}{}

%% The headline is where a themed deck puts its navigation bar. Emptying it entirely
%% leaves the frame title flush against the paper edge, so it carries the top margin
%% instead. This is the reliable place for that space: a \vspace* at the start of the
%% frametitle beamercolorbox lands *below* the title, not above it.
\setbeamertemplate{headline}{\vspace*{0.45cm}}

%% Page number only, muted, bottom right; nothing on the title frame.
\setbeamertemplate{footline}{%
  \ifnum\value{framenumber}>1\relax
    \hfill\usebeamerfont{footline}\usebeamercolor[fg]{footline}%
    \insertframenumber\hspace*{0.9cm}%
  \fi
  \vspace*{0.9ex}%
}

%% ---- frame title + short accent rule -----------------------------------------------
%% \rule rather than TikZ: no pgfpicture per frame, so compile time and PDF size stay
%% flat across ~35 frames.
%% No beamercolorbox: a full-width box has to re-derive the text margins by hand
%% (its leftskip/rightskip do not pick up \setbeamersize), which is what pushed the
%% title flush against the left edge. Typeset directly instead and beamer's own text
%% margins apply, so the title aligns with the body text on every frame. The trailing
%% negative skip claws back part of beamer's fixed frametitle-to-body separation,
%% which is tuned for themes whose title sits inside a filled band.
\setbeamertemplate{frametitle}{%
  \vspace*{0.30cm}%
  {\usebeamerfont{frametitle}\usebeamercolor[fg]{frametitle}%
   \raggedright\insertframetitle\par}%
  \ifx\insertframesubtitle\@empty\else
    \vspace*{0.35ex}%
    {\usebeamerfont{framesubtitle}\usebeamercolor[fg]{framesubtitle}%
     \raggedright\insertframesubtitle\par}%
  \fi
  \vspace*{0.9ex}%
  {\color{rpdaccent}\rule{2.4em}{1.1pt}}\par
  \vspace*{-0.10cm}%
}

%% ---- title page --------------------------------------------------------------------
%% Flush left, no band, no logo. Quarto's before-body.tex emits \frame{\titlepage}
%% automatically whenever `title` is set, so this template is the whole story.
%%
%% The venue lives here rather than in the qmd front matter because the title page can
%% only print the five variables beamer knows about (title, subtitle, author, institute,
%% date); an invented `venue:` key in the YAML is dropped in silence, exactly as the
%% deck's existing `email:` key already is. \providecommand, not \newcommand, so the
%% deck still compiles if the string is ever supplied from somewhere else. Blank it for
%% a non-conference showing and the date prints on its own -- see the guard below.
\providecommand{\rpdvenue}{ITAM Alumni Conference, Mexico City}
\setbeamertemplate{title page}{%
  \vfill
  {\usebeamerfont{title}\usebeamercolor[fg]{title}\raggedright\inserttitle\par}%
  \vspace{0.7em}{\color{rpdaccent}\rule{3.4em}{1.2pt}}\par
  \vspace{0.8em}%
  {\usebeamerfont{subtitle}\usebeamercolor[fg]{subtitle}\raggedright\insertsubtitle\par}%
  \vspace{1.8em}%
  {\usebeamerfont{author}\usebeamercolor[fg]{author}\insertauthor\par}%
  \vspace{0.15em}%
  {\usebeamerfont{institute}\usebeamercolor[fg]{institute}\insertinstitute\par}%
  \vspace{0.8em}%
  %% Venue and date share one line, and so share the `date' element's muted colour and
  %% \small sans face: no separate beamercolor/beamerfont is wanted here. The guard
  %% keeps an empty \rpdvenue from leaving a dangling separator in front of the date.
  {\usebeamerfont{date}\usebeamercolor[fg]{date}%
   \ifx\rpdvenue\@empty\else\rpdvenue{}\ \textperiodcentered\ \fi\insertdate\par}%
  \vfill
}

%% ---- section dividers --------------------------------------------------------------
%% Quarto's template already installs
%%   \AtBeginSection{\ifbibliography\else\frame{\sectionpage}\fi}      (l.57-63)
%% \AtBeginSection *replaces* rather than appends (beamerbasesection.sty:233,
%% \def\beamer@atbeginsection{#2}), so re-issuing it here swaps in a plain, unnumbered
%% divider instead of producing two. The \ifbibliography guard must be kept, or the
%% references section gets a divider of its own.
\AtBeginSection{\ifbibliography\else
  \begin{frame}[plain,noframenumbering]\sectionpage\end{frame}\fi}
\AtBeginSubsection{}
\setbeamertemplate{section page}{%
  \vfill
  {\usebeamerfont{section page}\usebeamercolor[fg]{section page}%
   \raggedright\insertsectionhead\par}%
  \vspace{0.6em}{\color{rpdaccent}\rule{3.4em}{1.2pt}}\par
  \vfill
}

%% ---- lists, blocks, captions -------------------------------------------------------
\setbeamertemplate{itemize item}{\raisebox{0.18ex}{\scriptsize\textbullet}}
\setbeamertemplate{itemize subitem}{\raisebox{0.20ex}{\tiny\textbullet}}
\setbeamertemplate{itemize subsubitem}{--}
\setbeamertemplate{enumerate item}{\insertenumlabel.}
\setbeamertemplate{blocks}[default]        % flat: no shadow, no rounded corners
\setbeamertemplate{bibliography item}{}

%% Slides are not a paper: drop the "Figure 1:" / "Table 1:" label and keep captions
%% as quiet notes. This has to go through \captionsetup rather than
%% \setbeamertemplate{caption}: Quarto's beamer template loads the caption package,
%% which takes over \caption entirely and leaves beamer's caption template unused.
%% Consequence: do NOT use @fig-/@tbl- cross references anywhere in the deck.
\captionsetup{
  labelformat = empty,
  labelsep    = none,
  font        = {scriptsize, color = rpdmute},
  justification = raggedright,
  singlelinecheck = false,
  skip        = 3pt
}

%% ---- helpers -----------------------------------------------------------------------
%% Full-bleed image on a {.plain} frame without adjustbox: an overwide box centred in
%% \textwidth spills symmetrically into both margins.
\newcommand{\bleed}[1]{\makebox[\textwidth][c]{#1}}

%% appendixnumberbeamer is not installed. Freeze the counter instead, so the backup
%% frames do not inflate the number the audience sees on the last content slide.
\newcounter{rpdmainframes}
\newcommand{\backupstart}{\setcounter{rpdmainframes}{\value{framenumber}}}
\newcommand{\backupend}{\setcounter{framenumber}{\value{rpdmainframes}}}

\makeatother
